\documentclass[10pt,a4paper]{amsart}
\usepackage[margin=19mm]{geometry}
\usepackage{amsmath,amssymb,mathtools}
\usepackage[scr=rsfs,cal=euler]{mathalfa}
\usepackage{xfrac}
\usepackage[unicode,hidelinks,bookmarks=false]{hyperref}
\newcommand{\ie}{i.e.\ }
\newcommand{\cf}{cf.\ }
\newcommand{\resp}{respectively\ }
\newcommand{\Subsection}{\subsection}
\providecommand{\nobreakdash}{\nobreak\hbox{-}}
\setlength{\emergencystretch}{2em}
\multlinegap0pt
\usepackage{amssymb}
\usepackage{enumerate}
\usepackage{float}
\usepackage{tikz}
\usepackage{bm}
\usepackage{csquotes}
\newtheorem{lemma}{Lemma}
\newtheorem{assumption}{Assumption}
\DeclareMathOperator{\Aut}{Aut}
\DeclareMathOperator{\Fr}{Fr}
\DeclareMathOperator{\Gal}{Gal}
\newcommand{\Q}{\numberset{Q}}
\newcommand{\T}{\boldsymbol{\mathrm{T}}}
\DeclareMathOperator{\Tr}{Tr}
\newcommand{\Z}{\numberset{Z}}
\newcommand{\calR}{\mathcal{R}}
\newcommand{\ga}{\alpha}
\newcommand{\gs}{\sigma}
\newcommand{\numberset}{\mathbb}
\title{Quaternionic Kolyvagin systems and Iwasawa theory for Hida families}
\author{Francesco Zerman}
\date{Source: arXiv:2507.04980v2 (CC BY 4.0)}
\begin{document}
\maketitle

\noindent\emph{Source and licence.} Quaternionic Kolyvagin systems and Iwasawa theory for Hida families. Version arXiv:2507.04980v2 (CC BY 4.0), distributed by arXiv under the Creative Commons Attribution 4.0 International licence, http://creativecommons.org/licenses/by/4.0/, as recorded in the arXiv licence metadata for this exact version and shown on the abstract page https://arxiv.org/abs/2507.04980v2. Full licence notice: \texttt{licenses/arxiv-2507.04980-CC-BY-4.0.txt}.\par\smallskip

\noindent\textit{Editorial scope.} Counted unit: the article's lemma lem:finer-choice-of-primes (source0.tex lines 1124--1126). The article's complete original proof of it is delivered with it (source0.tex lines 1127--1155, one \texttt{proof} environment of 29 lines). No mathematics is written, repaired or re-derived: the statement and the proof are the article's own lines, in the article's own notation, and no retained line is substituted or re-flowed.  Retained uncounted as the source-local context the proof needs: lines 641--647; lines 1099--1101.  Nothing was omitted from any retained span, so the package carries no omission record.  No retained line contains a citation, so the wrapper carries no bibliography and no bibliography entry.  No retained line contains a figure environment, an image-inclusion command or drawing code, so no image is delivered and the package declares no asset dependency.  Editorial formatting only, in addition: an AMS \texttt{amsart} preamble that restates the article's own declarations and macros used by the retained lines, namely the environments \texttt{lemma}, \texttt{assumption} on separate counters, together with the standard packages those lines need.  The retained environments print in this document as Lemma 1, Assumption 1 in order of appearance, whereas the article prints the same items with the numbering of its own full document.  The complete native source of the article as distributed in the arXiv e-print for this exact version is bundled unchanged as \texttt{sources/181141/source0.tex}.
\begin{lemma}
 \label{lem:charachteristic-polynomial-frobenius-Tm}
     Let $s\ge m$ and $\ell\in\mathcal{P}_{m,s}$. Then, $\Fr_\ell$ acts on $\T_m$ with characteristic polynomial
     \begin{equation*}
         X^2-(-1)^{\frac{k+j}{2}-1} T_\ell X+\ell.
     \end{equation*}
 \end{lemma}

\begin{assumption}\label{ass:big-image}
    The image of $G_\Q$ in $\Aut_{\mathcal{R}}(\T^{\dag})$ contains the scalars $\Z_p^\times$.
\end{assumption}

\begin{lemma}\label{lem:finer-choice-of-primes}
    For every $m,s\in\Z_{>0}$ with $s\ge m$ the set $\mathcal{L}_{m,s}$ is infinite.
\end{lemma}

\begin{proof}
    First, notice that an element $r\in\mathcal{R}_m$ is a unit if and only if its projection to one (hence any) nonzero quotient of $\calR_m$ is a unit.

    Fix a prime $\ell_0\in\mathcal{P}_{m,s}$. By Lemma \ref{lem:charachteristic-polynomial-frobenius-Tm}, we have that
    \begin{equation*}
        \Tr(\Fr_{\ell_0}\mid \T_m)=(-1)^{\frac{k+j}{2}-1}T_{\ell_0}\quad\text{and}\quad\det(\Fr_{\ell_0}\mid \T_m)=\ell_0.
    \end{equation*}
    Let now $\ga\in 1+p^s\Z_p$ and $\gs_\ga\in G_\Q$ such that its image in $\Aut(\T_m)$ is the scalar $\ga$, whose existence is granted by Assumption \ref{ass:big-image}. Then one obtains
    \begin{equation*}
        \Tr(\Fr_{\ell_0}\gs_\ga\mid \T_m)=(-1)^{\frac{k+j}{2}-1}\ga T_{\ell_0}\quad\text{and}\quad\det(\Fr_{\ell_0}\gs_\ga\mid \T_m)=\ga^2\ell_0.
    \end{equation*}
    Let $\mathcal{L}_{m,s}(\ell_0, \alpha)$ be the set consisting of those primes $\ell$ such that $\Fr_\ell$ is conjugated to $\Fr_{\ell_0}\gs_\ga$ in $\Gal(K(\T_{m,2s})/\Q)$. By Chebotarev's density theorem, the set $\mathcal{L}_{m,s}(\ell_0, \alpha)$ has infinite cardinality. For any $\ell\in\mathcal{L}_{m,s}(\ell_0, \alpha)$, one has
    \begin{equation}\label{eq:proof-finer-set-of-primes}
        T_{\ell}\equiv \ga T_{\ell_0}\quad\text{and}\quad \ell\equiv\ga^2\ell_0
    \end{equation}
    as elements of $\calR_{m,2s}$. Reducing to $\calR_{m,s}$ the above congruences and using the fact that $\ga\equiv 1\bmod p^s$, one easily sees that $\mathcal{L}_{m,s}(\ell_0, \alpha)\subseteq\mathcal{P}_{m,s}$. 

    We now look for an $\ga\in 1+p^s\Z_p$ such that the division of the two elements
    \begin{equation}\label{eq:elements-to-be-divided}
        \ga^2\ell_0+1\pm\ga T_{\ell_0}
    \end{equation}
    by $p^s$ is a unit of $\mathcal{R}_{m,2s}$. The existence of such an $\ga$ would imply that the set of primes $\mathcal{L}_{m,s}(\ell_0,\ga)\subseteq\mathcal{L}_{m,s}$, since, by \eqref{eq:proof-finer-set-of-primes}, we have the congruences
    \begin{equation*}
        \ell+1\pm T_\ell\equiv \ga^2\ell_0+1\pm\ga T_{\ell_0}
    \end{equation*}
    in $\calR_{m,2s}$, and this would conclude the proof.

    Writing $\alpha=1+p^{s}x$ for $x\in\Z_p$, dividing the elements \eqref{eq:elements-to-be-divided} by $p^s$ and reducing them modulo the maximal ideal of $\mathcal{R}_m$, one finds that in order to get such an $\alpha$ it is sufficient to require that the reduction of $x$ avoids two explicit values. Since $p>2$, this is always possible.
\end{proof}

\end{document}
